\chapter[Sermon 12. The Epistle of Thyatirena Is Expounded, Wherein Are Sundry Virtues Commended, and the Vice of Jezebel Reprehended.]{The Epistle of Thyatirena Is Expounded, Wherein Are Sundry Virtues Commended, and the Vice of Jezebel Reprehended.}
\label{sermon:012}
\markright{Sermon 12. The Epistle of Thyatirena Is Expounded, Wherein Are ...}

And unto the messenger of the congregation of Thyatira write. This says the Son of God, who has eyes like a flame of fire, and whose feet are like brass. I know your works and your love, service, and faith, and your patience, and your deeds, which are more at the last than at the first. Nevertheless, I have a few things against you, that you allow Jezebel, who calls herself a prophetess, to teach and deceive my servants, to lead them into fornication, and to eat foods offered up to idols.

\index{Repentance}The fourth epistle written to the Thyatireans, is more plentiful than the rest, and with manifold fruits replenished. For it commends and praises in that church excellent virtues, and singular gifts not a few. Straightforwardly he reproves them for suffering too gently the Jezebelism, which he describes and shows how filthy it is. He threatens them severely, unless with perfect repentance, they amend their sins and wickedness. Furthermore, he warns them that they should expect no new revelations, but that they persevere and abide in those which they had learned hitherto, and in which they now are. Hither also with most large promises he allures them, and finally communicates and commends this doctrine to all churches. And there is a wonderful likeness and correspondence in all epistles, as may be seen also in all the books of the prophets, in the story of the evangelists, and in Paul’s epistles. Whereof it may easily be gathered that the doctrine of the vexation is most absolute, perfect, and plain, and agreeable to itself in all things. In so much that if all the writings of all other Apostles and Prophets did remain, we should have had no more in those many and most plentiful books than we now have in the holy Bible. God provided well for us and for our infirmity by this brief way. Here are seven Epistles set in the 2nd Chapter: but it is marvelous to see how like they be all, teaching in a manner all one thing.

This fourth [section] is chiefly profitable for those congregations which are sound in the purity of doctrine and are pure moreover in holiness of life, but do not with a fervent zeal enough persecute open heresies. There are other fruits and benefits, which we shall speak of in order. But just as in all the other epistles that go before, first is set forth to whom the epistle is sent and from whom it comes: so also in this epistle to Thyatira, both the superscription, as they term it, and the subscriptio are expressly set. It is sent to the messenger of the church of Thyatira, and so to the whole church, as I have told you before often. And Thyatira is a noble and famous city of Lydia, in Asia, on the river Hermus: where we read that the woman was born, who sold purple, which was converted to Christ by Saint Paul in the 16th chapter of the Acts. It was a populous city and much frequented, so that it is no marvel though many diverse, unclean, curious, and heretics did associate and join themselves to the church of God. Geographers write many things of that famous city of Asia.

And the author of the epistle is the Lord Christ himself, the high king and Bishop, who uses the Apostle’s pen, or blessed John for his scribe or secretary, by whom he will have those things published throughout the whole world. And he establishes the epistle’s authority, while repeating certain members of the former image and description, he shows himself in such a way to be seen of the church, to be viewed in faith, that they help the matter wonderfully. He sees here heresies and the secrets of hearts, and treads under his most pure and clean feet whatever advances itself against God’s glory and truth.

He calls himself therefore the Son of God, whom before we heard to be the Son of Man. He is therefore and remains both, even in glory, as well the Son of God as man. In the divine nature, of the same substance with the Father, in the human nature, communicating with us in all things except sin, the other nature is not swallowed up in glory, but two distinct and several natures without any mixture abide in one person undivided: which indeed is one Christ, very God and very man, to be worshipped world without end. Hereof we have testimonies in Luke 1, 1 John 1, and Romans 1. And which of the heretics or persecutors will make war with the living Son of God?

After he attributes to himself eyes, casting out fire and flame. For nothing escapes the knowledge and judgment of Christ our Judge; he beholds the reins and hearts. Moreover, he lights some, and some he commits to everlasting fire, therein to burn forever. Now then, if any do imagine with themselves that they can hide heresies and malice in their hearts, they are deceived. For in the eyes of Christ, darkness itself is light also. The same Lord has also fetched most purged and clean ones; he treads down all ungodliness. And wherever he walks with his shining feet of brass, he consumes immediately all heresies and corrupt life. Therefore, this prelate, most pure, and most fit and apt to purge, finally best furnished to bolt out the secrets of hearts, shows to the congregations these things that follow: he himself walks and is conversant in the midst of the church, both King and Priest.

And just as he has testified in all his letters, that he knew the works of that church, so he repeats it here as well, so that we should never fall into wicked complacency, as though the almighty and all-knowing God did not know us and all that belongs to us—a matter about which I have spoken sufficiently before.

\index{John, Saint}\index{Arethas of Caesarea}Now he clearly sets forth every work of this congregation and commends five most notable gifts or brightest virtues. First, Charity, which encompasses the love of God and our neighbor: whereby it is brought to pass that we prefer nothing in the world before God, neither harm our neighbor, but rather heap upon him all duties and benefits. This we owe to God and all our brethren in the congregation. Of Charity is spoken elsewhere most abundantly, as in the gospel and epistle of Saint John. Secondly, he praises Diaconian, that is, the Ministry. The which may be explained two ways. For either he understands, as Arethas supposes, ministries toward the poor and needy, that is to wit, duties and pains taken about the poor, by leading, relieving, succoring, speaking faithfully in their cause, in giving them meat, drink, clothing, and visiting them. For so this word Diaconia is used in the second epistle to the Corinthians, and so forth. Or else he means the ministry of the word, by which in teaching, exhorting, comforting, and rebuking, we advance greatly God’s glory and the health of souls. The Thyatirenians were doubtless diligent in either of both. And accuse us grievously, which addict ourselves to our own affairs, do neglect our poor brethren: who finally make the ministry of God’s word odious, by our railing and slandering, especially with them that be ignorant as yet, and have heard nothing of God’s word.

\index{Aquinas, Thomas}He also commends faith among the Thyatireans. Thomas Aquinas in his commentary on this book admonishes that faith does not come from charity, because it is presented here first; but that charity and good works spring from faith. John has recited charity before faith, because faith derives its value from charity and works. Nevertheless, it seems here that “faith” is not so much to be taken as trust in God, but rather as fidelity, truth, and promises kept. For faithfulness beautifies all other gifts. Suppose you have servants, male and female, who are fortunate in doing their tasks, but are at the same time untrustworthy, fickle, and deceitful: what benefit would there be in their being equipped with various gifts? Imagine again that a preacher or senator is not so furnished with wisdom and experience, but is nevertheless faithful, and with all his heart strives to do all things uprightly and to favor the just cause: will not fidelity here supply his lack? Great therefore is faith, that is, fidelity and truth. Not without cause the Apostle required this of ministers in the 14th chapter of the first Epistle to the Corinthians, saying: That which is chiefly required of stewards is that a man be found trustworthy. This faith is also required of us at this day. This faith, good brethren, is rare, and therefore evils have overflowed everywhere. Let us heartily pray to the Lord that he will grant it to us, and that we may expel from our hearts unfaithfulness and deceitfulness.

Hereunto is added patience, which was also praised in the earlier churches. This is a necessary virtue. For impatience causes us to murmur and complain against God, and prevents us from standing firm in the confession of faith, while we refuse to suffer patiently the things that the enemies of faith threaten to inflict upon us. But why do you defile yourself with theft? Why do you rush into the wars of a foreign prince? Why do you practice usury and lewdness? Because you lack patience in your poverty, which you seek to relieve with wicked deeds.

To be brief, the Lord now recounts all kinds of good works. In this, he chiefly praises that they often surpassed themselves, in doing more and greater things. And this is worthy praise. For the husbandman, that is to say, the faithful father, prunes and cuts the vines, so that they may bear more plentiful fruit. It does not become the godly to remain stagnant, but to proceed in godliness.

It is most shameful to become worse the longer one endures. As a finger, the longer it is, the less it is. This is what is pointed out to children in schools who learn nothing. Let us be ashamed of our sluggishness. Let me say, weigh these things carefully in our minds, and consider often that God allows them, requires them, and that they are the true seals of the faithful walking in truth, and of those who boast of faith only as a vain name without the substance. If you feel yourself not to be utterly void of these gifts, praise God, and know that none of all these things is of yourself, but of grace. And pray for the increase of these gifts. If you are destitute of these virtues, mourn and lament before the Lord, humbly ask him for forgiveness, and require the abundance of God’s gifts.

Secondly, he rebukes certain things within that same congregation, namely that they allowed Jezebel to teach, and so forth. He calls this a minor matter, not that Jezebel’s doctrine is inherently insignificant, but because, though it is found among others, the church tolerated it gently—that is, did not punish it with greater severity. Regarding this phrasing, I have spoken of it previously. We do not condone Jezebel’s shameful acts, nor do we approve of them. But when we might have imposed a more severe punishment, we permitted them to abound and increase. Although, therefore, we possess many good gifts, the Lord is displeased with us for allowing ungodliness to reign.

But if the Lord should condemn that very allowance, how much more blameworthy should we consider the wickedness itself to be, namely, the Jezebelic practice? And how base and filthy it is, I will briefly declare.

Just as he previously refuted the Nicolaitans using the example of Balaam from Scripture, now he brings forth the example of Jezebel to refute the Cataphrygians, or Montanists. Arethas understands the entire passage concerning the Nicolaitans, which I do not dare to agree with, given the whole composition of the epistle. I grant that the Montanists were partakers in filthiness with the Nicolaitans. But Jezebel possesses a peculiar characteristic of her own.

Iezabel, as sacred history testifies in the third and fourth books of Kings, chapters 16 and 17 and following, was the daughter of Hethbahal, king of Sidon, who, when married to Ahab, introduced the worship of Baal into the kingdom of Israel, building a magnificent temple in Samaria and establishing a large college of Baal’s priests. For Elijah is recorded to have slain 450 Baalites, even of the kings’ chaplains, as it were canons or prebendaries, and 400 ministers or country chaplains who served in hills, woods, and groves. This woman, therefore, founded this religion and sought to govern the prophesying at her pleasure. For pursuing Elijah intensely, she slew many of the prophets, truly because they would not teach according to the woman’s desire. Moreover, through Baal’s religion, whoredom and all uncleanness increased. King Jehu objects to King Joram, his son, the whoredoms of his mother. So Iezabel also augmented the eating of meats offered unto idols and all idolatry throughout the whole kingdom. Even at that time, when the Lord in a solemn sacrifice by miracle on Mount Carmel through the ministry of Elijah, had declared to that whole realm that the religion of Baal was most vain and false, and that the religion of the only God of Israel was most sincere and true; for Iezabel nevertheless persecuted the truth and established falsehood. Yea, moreover, she took upon herself governance in civil matters.

She seized the king’s seal and forged letters, sending them in the king’s name to have Naboth put to death, a truly good and innocent man. Such was Jezebel’s wickedness.

\index{Apocalypse (Book of)}\index{Eusebius of Caesarea}\index{Epiphanius of Salamis}Now following the example of that defiled woman, there were women in the church of Thyatira who claimed authority in religion and teaching within the congregation, presuming to have the spirit of prophecy. They taught in fact, but corrupt doctrine, seducing those whom God by his doctrine had prepared to be his servants. But these false prophetesses corrupted their minds, and brought forth a new doctrine and prophecy, and many things not found in the scriptures, but drawn from their own devilish dreams and deceitfulness. And among other things, they communed with the Nicolaitans, in word, and participated in meals offered to idols, as has been spoken of before. And the Lord seems plainly to speak of the Cataphrygians or Montanists, whose foundation was laid in the time of Saint John, but after, in the course of time, and especially in the empire of Antoninus, sixty years after the Apocalypse was published, broke out more strongly and abundantly. They say how Montanus had prophetesses Priscilla and Maximilla, who had visions and brought in wonderful revelations into the church. Eusebius treats of them at length in the fifth book of ecclesiastical history, chapter 16. And Epiphanius in the 48th heresy, in Panarion. Certainly, John, or Christ himself by John, seeking at the very beginning to uproot and destroy the roots of this heresy, by the example of that wicked woman Jezebel, has condemned that same heresy. Scripture also elsewhere prohibits a woman from ruling, teaching, or ministering in the congregation. Immediately, the Lord himself will refute the new prophecies, when he admonishes us that he will reveal no other new kind of doctrine besides that which he has committed or delivered to his church. Now also fornication, and the eating of meats offered to idols, are condemned elsewhere in scripture most severely, as before is said.

But since those things so greatly afflicted and troubled the church of God in the time of the Apostles, it is not difficult to see how rash they are who, at this day (as I showed you before), for the sake of hating the true religion restored, accuse it of sects, which abound so plentifully, as though that corruption proved that the gospel we preach were not the gospel. For the gospel that was preached by John and the rest of the Apostles was the truest and purest gospel, even though false gospel proponents like the Nicolaitans, Cataphrygians, and other sects innumerable crept in. Nevertheless, the gospel refutes and condemns all such sects, and maintains the Christian truth and unity of the catholic church. Praise be to our God, the Lord. Amen.
